\section{引言}

无机固态电解质的离子电导率由局域配位势阱、载流子数量、可达空位、迁移网络连通性和骨架响应共同决定。静态晶体结构能够提供其中相当一部分几何信息，因此基于 CIF 的结构描述符常被用于高通量候选筛选\cite{Armand2008,Bachman2016,Adams2006,Sendek2017}。然而，跨 NASICON、硫化物和卤化物等体系直接汇总数据时，描述符分布和电导率分布都可能随材料类别整体移动。一个在全样本上相关性很高的特征，可能只是在识别“样本属于哪一类”，而不是捕捉可跨体系外推的离子传输规律。

除统计混杂外，结构描述符管线还面临更基础的软件风险。带氧化态的 \code{Na+}、\code{O2-} 等 Pymatgen \code{Species} 若被当作普通字符串比较，会漏检 Na、阴离子和部分占位；相对 CIF 路径若依赖当前工作目录，会在不同启动位置得到不同结果；全缺失特征若被静默插补，仍可能生成形式完整但没有结构信息的报告；在交叉验证之前对全数据插补和标准化，则会把测试折信息泄漏到训练过程。类似问题不会简单表现为程序崩溃，反而可能输出看似合理的数值，因此比显式错误更危险。

本文所描述的 \code{automat-naconductor} 修复分支，将研究问题重新定义为一个可审计的软件与统计契约：当结构输入不完整时，系统必须在创建研究结果之前失败；当输入可用时，每个筛选、组合和验证步骤都必须保留数据来源、控制变量、公式结构、缺失策略、随机种子和不可用原因。该框架不把相关性称为因果效应，也不把被跳过的验证策略当作零分。其目标不是在 84 个样本上得到尽可能高的训练性能，而是形成一个能够被逐项复核、可在补齐 CIF 后重新运行的结构描述符研究路径。

\statusbox{\textbf{当前执行状态。} 修复提交没有修改原始 CSV、任一 CIF、既有特征化文件或历史结果目录。当前 checkout 中原始表引用的 84 个 CIF 均不可用；Pipeline 和 Agent 都会在写入任何结果之前明确失败。因此，下文所有“候选”“验证”和“证据”均指算法设计与合成测试覆盖，不代表已经在真实 84 样本上重新获得材料学发现。}

\section{结果：从宽松原型到失败关闭的双轨研究框架}

\subsection{冻结输入、输出隔离与延迟写入}

新框架由一份可验证的 \code{run\_info.yaml} 作为单一运行契约。Pipeline 与 Agent 只共享三个冻结身份：原始 CSV 的规范路径、活动描述符注册表路径和注册表修订号。任何 Agent TSV 或审计 CSV 都必须包含并保持这三项身份唯一一致；混合批次、外来注册表或缺失身份均被拒绝。Pipeline 只能写入 \code{results/pipeline/}，Agent 只能写入 \code{results/agent/}，绝对路径、\code{..} 路径和跨轨目录均在 CLI 边界被拦截。

图~\ref{fig:dualtrack} 展示了两条轨道的关系。Pipeline 执行批量描述符审计、单描述符去混杂筛选、子采样 Lasso、物理族代表选择、受约束的公式搜索和 V1--V4 探索性验证；Agent 每次只评估一个明确指定的活动结构描述符，并使用与 Pipeline 相同的去混杂与折内 Ridge 交叉验证。两者可以并发启动，但在用户明确授权的 C9 之前，任何一方都不能读取另一方结果。C9 只允许对已完成且冻结的输出进行只读比较，一致性只能提高优先复核价值，不能构成因果证明。

\begin{figure}[htbp]
\centering
\resizebox{\textwidth}{!}{\begin{tikzpicture}[
  font=\sffamily\scriptsize,
  node distance=7mm and 8mm,
  box/.style={draw=NatureDark, rounded corners=2pt, align=center, minimum height=10mm, text width=31mm, fill=white},
  shared/.style={box, fill=NaturePale, line width=0.8pt},
  pipeline/.style={box, draw=NatureBlue, fill=NatureBlue!6},
  agent/.style={box, draw=NatureGreen, fill=NatureGreen!7},
  stop/.style={box, draw=NatureRed, fill=NatureRed!7, text width=33mm},
  arrow/.style={-{Latex[length=2mm]}, line width=0.7pt, draw=NatureGray},
  dasharrow/.style={-{Latex[length=2mm]}, line width=0.7pt, dashed, draw=NatureGray}
]

\node[shared, text width=104mm] (frozen) {冻结共享输入：raw CSV 规范路径 + 活动描述符注册表 + registry revision};
\node[shared, below=of frozen, text width=104mm] (preflight) {写前预检：CSV 列契约、全部 CIF 存在且可解析、结构描述符有有限值、输出路径不跨轨};
\node[stop, right=7mm of preflight, text width=25mm] (fail) {任一条件失败\\退出码 2\\不创建结果工件};

\node[pipeline, below left=10mm and 4mm of preflight] (p1) {Pipeline\\Stage 0--4};
\node[pipeline, below=of p1] (p2) {去混杂筛选\\子采样 Lasso};
\node[pipeline, below=of p2] (p3) {二/三元公式\\V1--V4 + CV};
\node[pipeline, below=of p3] (pout) {只写\\results/pipeline/};

\node[agent, below right=10mm and 4mm of preflight] (a1) {Agent\\显式单描述符};
\node[agent, below=of a1] (a2) {同一秩感知控制\\同一折内 Ridge CV};
\node[agent, below=of a2] (a3) {TSV 指标 +\\结构审计 CSV};
\node[agent, below=of a3] (aout) {只写\\results/agent/};

\node[shared, below=13mm of $(pout)!0.5!(aout)$, text width=80mm] (c9) {C9（需用户授权）：两轨均完成并冻结后，只读比较；一致仅为三角验证线索，causal claim = false};

\draw[arrow] (frozen) -- (preflight);
\draw[arrow] (preflight) -- (p1);
\draw[arrow] (preflight) -- (a1);
\draw[arrow] (p1) -- (p2);
\draw[arrow] (p2) -- (p3);
\draw[arrow] (p3) -- (pout);
\draw[arrow] (a1) -- (a2);
\draw[arrow] (a2) -- (a3);
\draw[arrow] (a3) -- (aout);
\draw[arrow] (pout) -- (c9);
\draw[arrow] (aout) -- (c9);
\draw[dasharrow] (preflight.east) -- (fail.west);

\draw[decorate,decoration={brace,amplitude=4pt},NatureGray] ($(p1.north west)+(-2mm,2mm)$) -- ($(pout.south west)+(-2mm,-2mm)$) node[midway,left=5mm,align=center]{批量\\组合轨};
\draw[decorate,decoration={brace,amplitude=4pt,mirror},NatureGray] ($(a1.north east)+(2mm,2mm)$) -- ($(aout.south east)+(2mm,-2mm)$) node[midway,right=5mm,align=center]{独立\\审计轨};
\end{tikzpicture}
}
\caption{\textbf{失败关闭的双轨结构描述符框架。} 冻结输入首先经过路径、CIF 可解析性和注册表身份预检。Pipeline 与 Agent 共享输入但隔离输出；任何缺失或身份冲突均在写工件前终止。只有在用户授权的 C9 中，才允许对已冻结结果作只读三角比较。}
\label{fig:dualtrack}
\end{figure}

这一设计改变了输出创建时机。特征化必须先解析相对路径并严格预检全部 CIF，Stage 0 必须确认存在足够有效的真实描述符，才创建 \code{results/pipeline/}。Agent 也必须在任何 TSV、审计 CSV 或图件写入之前完成所有 CIF 的存在性和可解析性检查。缺少 84 个 CIF 时，两个入口均以受控诊断和退出码 2 结束，而不是留下空目录、全 NaN 表格或 Python traceback。

\subsection{结构描述符注册表从函数清单升级为可搜索契约}

注册表保留 41 个公开 callable，以兼容旧代码，但新增单位、物理维度、活动状态和别名元数据。三个条目不再进入自动搜索：\code{max\_bond\_length} 是 \code{a2\_max\_dist} 的兼容别名；\code{bottleneck\_anisotropy} 当前没有有效实现；\code{bvse\_barrier\_estimate} 在没有 BVSE 后端时永久返回不可用。其余 38 个描述符构成活动搜索空间。与旧稿不同，\code{wyckoff\_diversity} 已通过 \code{SpacegroupAnalyzer.get\_symmetrized\_structure()} 的等价索引重新实现，不再被视为全缺失列。

\begin{table}[htbp]
\centering
\caption{八个物理族及当前活动搜索规模。公开数量包括兼容 callable；活动数量排除了别名和无实现占位符。}
\label{tab:families}
\begin{tabularx}{\textwidth}{@{}lXcc@{}}
\toprule
族 & 物理含义 & 公开数 & 活动数 \\
\midrule
A & Na 配位多面体：键长、配位数、体积、畸变与方向性 & 11 & 9 \\
B & Na--Na 网络：邻接、连通分量和网络维度 & 5 & 5 \\
C & Na 浓度与占位：数密度、占位和位点数 & 3 & 3 \\
D$'$ & 周期性空位拓扑：Voronoi 间隙、可达性与网络维度 & 5 & 4 \\
E & 骨架刚性与多面体共享 & 4 & 4 \\
F & 长程几何：次近邻、角度、均匀性与曲折度 & 4 & 4 \\
G & 电子化学代理，全部预标记为高风险 & 4 & 4 \\
H & 对称性和占位异质性，部分预标记为高风险 & 5 & 5 \\
\midrule
合计 &  & 41 & 38 \\
\bottomrule
\end{tabularx}
\end{table}

化学物种识别统一转换为元素符号并按元素累计占位，避免 \code{Na+} 与字符串 \code{Na} 不相等造成的漏检。G 族电荷平衡偏离在存在氧化态时改为
\begin{equation}
D_q=\frac{|\sum_i q_i|}{\sum_i |q_i|},
\end{equation}
只有缺少氧化态时才使用明确标记的回退估计。D$'$ 族不再把 Voronoi ridge 中心误作空隙点，而是从真实 Voronoi 顶点生成候选，并以晶格最小像计算间隙--Na 距离和周期去重。该修改避免跨周期边界约 $0.2$~\AA{} 的近邻被错误记录为接近一个晶胞长度。

\subsection{秩感知控制揭示当前分类设计的结构性共线}

当前 84 行表的分类分布为：30 个 NASICON 样本全部标记为 oxide，41 个 sulfide 样本全部标记为 sulfide，13 个 halide 中有 12 个 chloride 和 1 个 iodide。这里的“NASICON--oxide”或“sulfide--sulfide”冗余不是材料学恒等式，而是当前采样设计的线性代数事实。若体系类别已经以参考编码虚拟变量进入控制矩阵，oxide 和 sulfide 对比不再增加设计矩阵秩；唯一新增信息来自 halide 内的 iodide--chloride 对比。

新实现先构建带截距的 system 设计，再逐列测试 anion 对比是否增加矩阵秩（图~\ref{fig:rankcontrol}）。当前审计得到 system 设计秩为 3，加入全部 anion 列后的秩为 4，阴离子增量秩为 1。实际残差化只使用 system 主控制和能够增加秩的阴离子列，同时把冗余列、增量列和最终残差化列写入结果属性。\code{anion\_is\_independent\_control=False} 明确禁止把阴离子类别解释为独立于体系的完整控制因素。

\begin{figure}[htbp]
\centering
\resizebox{\textwidth}{!}{\begin{tikzpicture}[
  font=\sffamily\scriptsize,
  cell/.style={draw=NatureLight, minimum width=20mm, minimum height=8mm, align=center},
  head/.style={cell, fill=NaturePale, font=\sffamily\bfseries\scriptsize},
  rank/.style={draw=NatureBlue, rounded corners=2pt, fill=NatureBlue!6, align=left, text width=48mm, inner sep=3mm},
  arrow/.style={-{Latex[length=2mm]}, draw=NatureGray, line width=0.7pt}
]

\matrix (m) [matrix of nodes, nodes={cell}, row sep=-\pgflinewidth, column sep=-\pgflinewidth] {
 |[head]|system & |[head]|chloride & |[head]|iodide & |[head]|oxide & |[head]|sulfide \\
 NASICON (30) & 0 & 0 & 30 & 0 \\
 sulfide (41) & 0 & 0 & 0 & 41 \\
 halide (13) & 12 & 1 & 0 & 0 \\
};

\node[rank, right=10mm of m] (r1) {\textbf{步骤 1：system 主控制}\\带截距、参考编码\\设计秩 $=3$};
\node[rank, below=6mm of r1] (r2) {\textbf{步骤 2：逐列测试 anion 对比}\\oxide、sulfide：不增加秩\\iodide vs chloride：增加 1 个秩};
\node[rank, below=6mm of r2] (r3) {\textbf{实际残差化矩阵}\\system + 唯一增量 anion 列\\combined rank $=4$\\anion independent control = false};
\node[draw=NatureRed, rounded corners=2pt, fill=NatureRed!6, align=center, text width=48mm, below=6mm of r3] (skip) {iodide 仅 1 例\\阴离子分层 CV：\textbf{SKIPPED}\\不是 0 分，也不是失败证据};

\draw[arrow] (m.east) -- (r1.west);
\draw[arrow] (r1) -- (r2);
\draw[arrow] (r2) -- (r3);
\draw[arrow] (r3) -- (skip);
\end{tikzpicture}
}
\caption{\textbf{当前 84 行分类设计的秩感知控制。} system 先进入参考编码设计；oxide 和 sulfide 阴离子对比与 system 完全共线，只有 iodide 相对 chloride 的对比增加一个秩。阴离子分层交叉验证因 iodide 仅一例而必须显式跳过。}
\label{fig:rankcontrol}
\end{figure}

对每个活动描述符 $X_j$ 和目标 $Y=\log_{10}(\sigma/\mathrm{S\,cm^{-1}})$，分别用 Ridge 回归移除控制矩阵 $Z$ 能解释的部分：
\begin{align}
\widehat{X}_j &= f_{\mathrm{Ridge}}(Z), & X^{\mathrm{res}}_j &= X_j-\widehat{X}_j,\\
\widehat{Y} &= g_{\mathrm{Ridge}}(Z), & Y^{\mathrm{res}} &= Y-\widehat{Y},\\
\rhoDeconf,j &= \rho_s(X^{\mathrm{res}}_j,Y^{\mathrm{res}}).
\end{align}
Ridge 用于稳定小样本和共线控制下的残差估计\cite{Hoerl1970}。该值是控制已观测分类因素后的单调关联，不是处理效应或因果参数\cite{Pearl2009,Chernozhukov2018}。

体系代理比按实际代码处理边界：当原始相关平方低于 $10^{-12}$ 时设为 0；当原始相关与去混杂相关符号相反时设为 1；其余情况计算
\begin{equation}
P_j=\operatorname{clip}\left(1-\frac{\rhoDeconf,j^2}{\rhoRaw,j^2},0,1\right).
\end{equation}
先以 $|\rhoRaw|<0.2$ 标记噪声级；否则 $|\rhoDeconf|>0.3$ 标记强物理信号；剩余描述符按 $P_j<0.3$、$0.3\leq P_j<0.7$ 和 $P_j\geq0.7$ 分别标记为弱物理信号、混合信号和体系代理。这些阈值是候选压缩规则，不等同于多重比较校正后的显著性声明。

\subsection{所有统计预处理被移入训练折或子样本内部}

旧实现曾在全数据上完成中位数插补和标准化，再进入交叉验证与稳定性筛选。修复后，描述符原始值和 NaN 被原样保留。活动列必须至少覆盖 50\% 样本，低覆盖列被排除；随后加入 15 个由固定种子 42 生成的标准正态噪声列。真实值与噪声都不在全局阶段拟合插补器或缩放器。

稳定性选择改为真正的子采样 Lasso\cite{Tibshirani1996,Meinshausen2010}。每次从 $n$ 个样本中无放回抽取 $\lfloor0.5n\rfloor$ 个样本，并在该子样本内部拟合
\begin{equation}
\texttt{SimpleImputer(median)}\rightarrow\texttt{StandardScaler}\rightarrow
\texttt{Lasso}(\alpha=0.05).
\end{equation}
只有绝对值大于 $10^{-12}$ 的非零系数才计为一次选择。100 次子采样后，真实特征需要同时满足选择频率 $>0.60$ 和高于 15 个噪声列选择频率的第 95 百分位数。Stage 2 只接受 Stage 1 已通过的真实描述符和固定噪声，不能重新引入被去混杂筛选拒绝的特征。

族代表按 $|\rhoDeconf|$ 排序但保留方向。二元公式模式每族最多保留一个稳定代表；三元模式每族最多保留两个。第二个同族名额只为“同族两个加相邻族一个”的受限三元语法服务，不应被解释为第二项独立科学发现；缺失其它族时也不会把该硬上限扩张到更多名额。

\subsection{声明式公式语法替代无界代数枚举}

组合搜索直接在原始物理值上完成，只有完整公式进入模型后才在训练折内部插补和标准化。加法和乘法通过 \code{itertools.combinations} 只保留一个规范无序对；比例只允许规则表中明确声明的方向，当前自动搜索唯一比例方向为 A$\rightarrow$C。分母为零或任一分量非有限时，结果保持 NaN，不加入 epsilon，也不以极大值替代。

\begin{table}[htbp]
\centering
\caption{当前自动公式搜索的跨族运算规则。未列出的族对不进入自动枚举；同族 A--H 均允许加法和乘法。}
\label{tab:rules}
\begin{tabularx}{\textwidth}{@{}lllX@{}}
\toprule
第一族 & 第二族 & 允许运算 & 方向与限制 \\
\midrule
A & B & $+$、$\times$ & 交换运算，仅生成一个规范无序对 \\
A & D$'$ & $\times$ & 交换运算 \\
A & C & $/$ & 只允许 A/C，不自动生成 C/A \\
B & D$'$ & $+$、$\times$ & 交换运算 \\
A & H & $+$、$\times$ & 交换运算，高风险 H 仍保留 provenance \\
A & E & $\times$ & 交换运算 \\
\bottomrule
\end{tabularx}
\end{table}

加法只有在两个已知物理维度相同时允许；未知自定义描述符为了兼容外部注册表可继续进入审查。乘法和比例会记录组合后的维度表达。禁止 \code{log}、\code{sqrt} 和幂运算。三元公式只允许 2 或 3 个描述符，结构固定为同族两个描述符先以加法或乘法结合，再与规则表中显式相邻族的一个描述符结合；任意三元组不被枚举。每个候选必须携带 \code{components}、\code{operators}、\code{component\_families}、规则 ID、原始值来源、缺失策略和标准化时机。CSV 使用严格 JSON 保存这些容器，Stage 4 重新读取时会校验组分数、运算符数、\code{d1/d2/operator} 与 provenance 是否一致，任何损坏都直接报错，绝不把三元式静默降级为二元式。

\begin{figure}[htbp]
\centering
\resizebox{\textwidth}{!}{\begin{tikzpicture}[
  font=\sffamily\scriptsize,
  box/.style={draw=NatureDark, rounded corners=2pt, align=center, minimum height=10mm, fill=white},
  rule/.style={box, fill=NaturePale, text width=29mm},
  evidence/.style={box, text width=27mm, minimum height=22mm},
  arrow/.style={-{Latex[length=2mm]}, draw=NatureGray, line width=0.7pt}
]

\node[rule] (raw) {原始物理值\\NaN 原样保留};
\node[rule, right=4mm of raw] (pair) {二元规则\\同族 $+,\times$\\显式跨族表};
\node[rule, right=4mm of pair] (triple) {受限三元\\同族两个 +\\显式相邻族一个};
\node[rule, right=4mm of triple] (prov) {严格 provenance\\组分、运算符、维度\\规则 ID 与缺失策略};

\draw[arrow] (raw) -- (pair);
\draw[arrow] (pair) -- (triple);
\draw[arrow] (triple) -- (prov);

\node[evidence, draw=NatureBlue, fill=NatureBlue!6, below=12mm of raw] (v1) {\textbf{V1}\par 匹配噪声基线\par system 内逐分量置换\par 100 draws};
\node[evidence, draw=NatureGreen, fill=NatureGreen!7, right=4mm of v1] (v2) {\textbf{V2}\par 折安全目标残差预测\par rank-aware controls\par OOF Spearman};
\node[evidence, draw=NatureDark, fill=NaturePale, right=4mm of v2] (v3) {\textbf{V3}\par 每个 system 内\par raw Spearman\par 独立可用状态};
\node[evidence, draw=NatureRed, fill=NatureRed!6, right=4mm of v3] (v4) {\textbf{V4}\par system 分层 bootstrap\par 500 draws\par 95\% CI};

\draw[arrow] (prov.south) |- ($(v2.north)!0.5!(v3.north)$);
\draw[arrow] ($(v2.north)!0.5!(v3.north)$) -| (v1.north);
\draw[arrow] ($(v2.north)!0.5!(v3.north)$) -- (v2.north);
\draw[arrow] ($(v2.north)!0.5!(v3.north)$) -- (v3.north);
\draw[arrow] ($(v2.north)!0.5!(v3.north)$) -| (v4.north);

\node[box, below=8mm of $(v2.south)!0.5!(v3.south)$, text width=105mm, fill=NaturePale] (cv) {三种 CV 独立运行：anion-stratified（不足则 skip/降折）｜LOSO｜10 次 system 分层 80/20 子采样。综合分只平均有限且未 skip 的绝对 Spearman，并报告可用策略数。};
\draw[arrow] (v1.south) |- (cv.west);
\draw[arrow] (v2.south) -- (cv.north west);
\draw[arrow] (v3.south) -- (cv.north east);
\draw[arrow] (v4.south) |- (cv.east);
\end{tikzpicture}
}
\caption{\textbf{受约束公式语法与 V1--V4 探索性证据。} 候选只来自显式规则，完整公式在原始值上构造。V1 以体系内分量置换生成匹配噪声基线；V2 用折安全的已知因素残差预测；V3 报告各体系内部原始 Spearman；V4 给出体系分层 bootstrap 区间。所有区块均带可用状态和原因。}
\label{fig:formula}
\end{figure}

\subsection{V1--V4 与可独立跳过的交叉验证}

Top-$k$ 候选并不只按一个综合分验证，而是生成四个命名证据区块（图~\ref{fig:formula}）。

\begin{enumerate}
  \item \textbf{V1 匹配噪声基线。} 对公式的每个分量在各 system 内独立置换，再用相同运算符重建公式。默认 100 次置换，报告噪声绝对 Spearman 的中位数、第 95 百分位数和观测值所处分位。
  \item \textbf{V2 已知因素后的折安全预测。} 在训练折内用秩感知控制矩阵拟合目标并形成训练、测试残差，再用仅含公式值的折内中位数插补--标准化--Ridge 模型预测测试残差。主量为 held-out 残差与公式预测的 OOF Spearman；双侧偏相关只作为补充。该区块明确标记 \code{causal\_claim=false}。
  \item \textbf{V3 体系内关联。} 分别在每个 system 内计算原始 Spearman，样本不足或常数数据时单独标记不可用。
  \item \textbf{V4 体系分层 bootstrap。} 在每个 system 内有放回抽样，再拼接各体系样本；默认 500 次，报告原始公式--目标 Spearman 的 95\% 分位区间。该区间不包含筛选与公式选择的不确定性。
\end{enumerate}

常规预测诊断使用同一个折内 \code{Imputer--Scaler--Ridge} 管线。阴离子分层验证请求 3 折，但若任一类别少于 2 个样本则显式 \code{skipped}；若最小类别支持 2 而请求 3 折，则自动降为 2 折并报告 requested/effective folds。当前 iodide 只有一例，因此阴离子分层策略不可用。LOSO 需要至少两个 system；重复子采样严格执行 10 次、测试比例 0.2 和种子 42，并检查训练集和测试集能否覆盖所有 system。

综合分只平均有限且未跳过策略的绝对 Spearman：
\begin{equation}
S=\frac{1}{m}\sum_{k\in\mathcal A}|\rho_k|,
\end{equation}
其中 $\mathcal A$ 为实际可用策略集合，$m=|\mathcal A|$。输出同时记录可用策略数和 \code{composite\_is\_complete}。缺少一种验证不会被当作 0 分，但也不能被误写为“三种策略全部支持”。

\subsection{Agent 轨道提供单描述符独立审计而非第二个黑箱模型}

Agent 不再使用组成特征、随机森林或预拆分 train/validation/test 文件。每次运行必须显式指定活动描述符名，例如 \code{a2\_max\_dist}。它对全量冻结 raw CSV 执行全部 CIF 预检、计算单一结构描述符、调用同一个秩感知去混杂分析器和多策略折内 Ridge CV，并把材料标识、解析后的 CIF 路径、体系、阴离子、目标值和该描述符原始值写入审计 CSV。

Agent 状态判断只读取自己的 TSV。只有 \code{evaluated}、\code{keep} 或 \code{discard} 且去混杂 Spearman 有限的行才用于耐心判断；\code{crash} 和不可用 CV 策略不会被伪装成“没有改善”。现有/新 TSV 和审计 CSV 在覆盖前都必须通过冻结批次身份检查。该轨道的价值是从不同搜索机制审计单个结构变量，而不是复制 Pipeline 的组合选择。

\subsection{验证范围与当前不能下的结论}

修复分支在项目目录中通过 70 项自动化测试、Python 全量编译和五个 CLI 帮助入口检查。测试覆盖带电物种、混合占位、周期最小像、Wyckoff 多样性、全缺失失败关闭、折内插补、不可行 CV 的 skip 语义、Stage 1--2 合约、硬族容量、真实可达三元式、严格 JSON 往返、V1--V4 provenance、双轨输出隔离和冻结身份拒绝。

这些测试证明的是代码契约和合成数值路径按设计运行，不证明 38 个活动描述符中哪一个对真实材料最重要。由于当前 CIF 不可用，不能从该分支声称“最佳单描述符”“最佳组合”“组合提升百分比”或新的传输机制。补齐并冻结结构后，真实 Pipeline 应从 Stage 0 重新生成特征，而不能继续使用旧的全缺失特征化表。

\section{讨论}

该修复的关键变化是把“能输出一个结果”从成功标准中移除。材料信息学管线最危险的状态不是报错，而是在结构输入、分类设计或验证策略不满足前提时仍生成整齐的排名。失败关闭、显式 skip 和 provenance 使“没有证据”与“负证据”得到区分：CIF 缺失不是低性能；iodide 仅一例不是阴离子 CV 得分为零；LOSO 不可行也不应杀死其它策略。

秩感知控制还说明，混杂变量的名称并不能保证其统计独立性。当前数据中 system 与大多数 anion 列几乎完全重合，机械地同时 one-hot 编码会制造冗余控制并掩盖采样设计限制。先审计秩、再残差化，比在方法部分笼统写“控制体系和阴离子”更可核查。

受限公式搜索牺牲了搜索空间的充分性，但换来规则可读、方向明确和量纲可审计。当前规则表仍是一组人工先验，而不是材料学定律。特别是 A/C 比例方向、A--H 组合和三元邻接关系，在扩展前需要逐项物理审查。V1--V4 也仍是探索性证据，因为外层 group split 中没有重新执行“筛选--稳定性--代表--组合选择”的完整 nested reselection；V4 区间只反映固定公式在当前分层抽样下的波动。

最终，这一框架更适合作为候选复核与计算验证的前置层，而不是终点。通过结构数据审计、去混杂和跨域诊断获得的候选，仍需结合可追踪文献电导率、位点占据合理性、BVSE/BVEL 或几何通道分析，以及后续 DFT、NEB 或 AIMD 进行验证。本文没有运行这些计算，也不把静态 CIF 相关性替代为迁移势垒或有限温度动力学结论。
